\begin{korollar}{Existenz fast isometrischer Abbildungen}
                {ExistenzFastIsometrischerAbbildungen}
Seien $n \in \N$ und $(E, \norm\cdot_E)$ ein normierter, $n$--dimensionaler $\R$--Vektorraum.\\
Seien weiter $\eps\in\!\ ]0,1[$, und $m \in \N$ mit
$m \le \frac 2 \pi \cdot \left(\frac{\eps}
{d_{BM}((E, \norm\cdot_E), \ell_2^n)} \right)^2 n$.\\
Dann existiert ein injektives $A \in L(\ell_2^m, E)$ mit:
	\[\opnorm A{\ell_2^m}{E}\opnorm{A\inv}{A\left(\ell_2^m\right)}{\ell_2^m}
  \le \frac {1+\eps}{1-\eps}\text.\]
Insbesondere existiert $F \le E$ mit $\dim F = m$ und 
\[d_{BM}\left(\left(F,\restrict{\norm\cdot_E}F\right),\ell_2^m\right)
\le\frac{1+\eps}{1 - \eps}\text.\]
\end{korollar}
\begin{proof}
Sei $d \coloneqq  d_{BM}((E, \norm\cdot_E), \ell_2^n)$.\\
Nach~\ref{NormäquivalenzbasisImEndlichdimensionalen} existiert eine Basis 
$x \in E^n$ mit
$\frac 1 d \norm\cdot_{2, n} \le \norm\cdot_E \circ \bprod \cdot x_n
\le \norm\cdot_{2, n}$.\hfill $(\ast)$\\
Seien $T\coloneqq S^{m-1}$,
\[\bigg((\Omega, \mfa, \Prim), (\gamma,\delta,\widetilde g),(n, m, mn)\bigg)
\text{ ein }(N(0,1), 3)\text{--Tupel zu }(3,(n,m,mn))\] (existiert nach
\ref{ExistenzNormalNTupel}) und $(Z_{\gamma, x}, M_{\delta,T})$ das
Auswertungspaar zu $(\gamma, \delta, x, T)$.\\
Wegen $(\ast)$ gilt insbesondere $x \neq 0$ und
$\norm\cdot_E \circ \bprod\cdot x_n \le \norm\cdot_{2, n}$.\hfill $(a)$\\
Es gilt: $T \neq \emptyset$, $T = - T$ und $T \subseteq S^{m-1}$.\hfill $(b)$\\
Schließlich gilt:
\begin{align*}
\E\left(M_{\delta,T}\right)
&=\E\left(\sup_{t\in S^{m-1}}\bprod t\delta_m \right)
\overset{\ref{IterierteEuklidischeNorm}}=
\E\left(\norm\cdot_{2,m}\circ \delta \right)
\overset{\ref{ErwartungswertDer2NormEinesStandardnormalverteiltenProzesses}}
=a_m
\overset{\refclaim{AbschaetzungEinesGammafunktionsquotienten}
{AbschaetzungEinesGammafunktionsquotienten2}}\le \sqrt m
\overset{\text{Vor.}}\le \frac \eps d \sqrt{\frac 2 \pi n}\\
&\overset{\refclaim{AbschaetzungEinesGammafunktionsquotienten}
{AbschaetzungEinesGammafunktionsquotienten4}}\le  \frac \eps d a_n
\overset{\ref{ErwartungswertDer2NormEinesStandardnormalverteiltenProzesses}}=
\frac \eps d\E\left(\norm\cdot_{2,n} \circ \gamma \right)
\overset{(\ast)}\le 
\eps \E\left(\norm\cdot_E \circ \bprod \cdot x_n \circ \gamma \right)\\
&= \eps\E\left(Z_{\gamma,x} \right)\text.\tag {$c$}
\end{align*}
Wegen $(a), (b), (c)$ sind die Voraussetzungen von
\ref{ExistenzEingeschraenkterFastIsometrien} erfüllt und wegen $T = S^{m-1}$
existiert nach~\ref{ExistenzEingeschraenkterFastIsometrien} ein
$\emptyset \neq \mca \subseteq L(\ell_2^m, E)$, sodass jedes $A \in \mca$
injektiv ist und gilt:
\[\opnorm A {\ell_2^m}E \opnorm{A\inv}{A(\R^m)}{\ell_2^m} \le
\frac {1+\eps}{1-\eps}\text{.}\]
Wegen $\mca \neq \emptyset$ ist damit die Aussage (ohne den Zusatz) gezeigt.\\
Seien $A_0 \in \mca$ und $F\coloneqq  \Bild A_0$. Da $A_0$ injektiv ist, ist
$\dim F = m$ und damit ist $A_0 \in \Inv(\ell_2^m, F)$, sodass gilt:
\begin{align*}
d_{BM}\left(\ell_2^m, \left(F, \restrict{\norm\cdot_E}F\right) \right)
&= \inf\menge{\opnorm S {\ell_2^m}{F}\opnorm {S\inv}F{\ell_2^m}}
{S \in \Inv(\ell_2^m, F)}\\
& \le \opnorm {A_0} {\ell_2^m}E \opnorm{{A_0}\inv}{F}{\ell_2^m}
\le \frac {1+\eps}{1-\eps}\text.
\end{align*}
%Für alle $\omega \in \Omega$ gilt dann:
%5\begin{align*}
%(\norm\cdot_E \circ \bprod \cdot x_N\circ \gamma )(\omega)
%&= \norm{\bprod{\gamma(\omega)}x_N}_E \\
%&\overset{(\ast)}\ge 
%\frac 1 d \norm{\gamma(\omega)}_{2, N} = 
%\left(\frac 1 d \norm\cdot_{2, N} \circ \gamma \right)(\omega) \quad (\ast\ast)
%\end{align*}
%und
%\begin{align*}\sup_{t \in S^{m-1}}\bprod t \delta_m
% \overset{IterierteEuklidischeNorm}= \norm\cdot_{2, m}\circ \delta\quad 
%(\ast\ast\ast)\end{align*}
%Die jeweils ersten Abbildungen in den (Un)Gleichungen sind nach
%\refclaim{IntegrierbarkeitVonExtremaInGordonschenPaaren}
%{IntegrierbarkeitVonExtremaInGordonschenPaaren1} integrierbar, die jeweils
%letztgenannten sind es nach
%\ref{ErwartungswertDer2NormEinesStandardnormalverteiltenProzesses}.\\
%Wegen $(\ast)$ ist insbesondere $\norm\cdot_E\circ \bprod \cdot x_N \le
%\norm\cdot_{2, N}$.\\
%Ferner gilt $S^{m-1}= -S^{m-1}$ und
%\begin{align*}
%\E\left(\sup_{t \in S^{m-1}}\bprod t \delta_m \right)
%&%\overset{
%\substack{\ref{ErwartungswertDer2NormEinesStandardnormalverteiltenProzesses},
%\\(\ast\ast\ast)}}=
%a_m \overset{\refclaim{AbschaetzungEinesGammafunktionsquotienten}
%{AbschaetzungEinesGammafunktionsquotienten2}}\le \sqrt m \le
%\sqrt {\frac 1 2}\sqrt N \frac{\eps}d 
% \overset{\refclaim{AbschaetzungEinesGammafunktionsquotienten}
%{AbschaetzungEinesGammafunktionsquotienten4}}\le \eps \frac 1 d a_N\\
%&
%\overset{\ref{ErwartungswertDer2NormEinesStandardnormalverteiltenProzesses}}= 
%\E\left(\eps \frac 1 d \norm\cdot_{2, N}\circ \gamma  \right)
%\overset{(\ast\ast)}\le \eps\E(\norm\cdot_E\circ \bprod \cdot x_N \circ \gamma)
%\end{align*}
%Damit sind die Voraussetzungen von~\ref{ExistenzBanachMazurnaherAbbildungen}
%erfüllt und es gibt ein $A \in L(\R^m, E)$ injektiv mit
%$\opnorm A {\R^m}E\opnorm{A\inv}{A(\R^m)} {\R^m} \le \frac {1+\eps}{1-\eps}$.\\
%Sei $F\coloneqq \Bild A$. Dann ist $F \le E$ und da $A$ injektiv ist, ist $\dim F=m$.\\
%Ferner gilt wegen der Surjektivität von $A$ auf $F$:
%\begin{align*}
%d_{BM}((F, \norm\cdot_E\vert_{F}), l_2^m) &= 
%\inf\menge{ \opnorm{B}{\R^m}F \opnorm {B \inv}F {\R^m}}{B \in \Inv(F, \R^m)}\\
%&\le \opnorm A {\R^m}E\opnorm{A\inv}{A(\R^m)} {\R^m} \le \frac {1+\eps}{1-\eps}
%\end{align*}
\end{proof}